\paragraph{An ordinary Jacobian in the real-multiplication class.}
The ordinary curve already used in A13,
\[
 C/\mathbf F_2:\qquad y^2+y=x+x^{-1}+(x-1)^{-1},
 \tag{RME.1}
\]
has a Rosati-invariant integral action of
$\mathcal O_5=\mathbf Z[(1+\sqrt5)/2]$ on its Jacobian $A$,
defined over $\mathbf F_2$. Thus the ordinary Jacobian class
appearing in the refined assertion of A33 is nonempty.

\begin{proof}
The Frobenius polynomial checked in A13 is
$P_2(X)=X^4-X^2+4$. Let $\pi$ be $2$-power Frobenius and
let $V=\pi^\dagger=2/\pi$ be its Verschiebung, using the
canonical principal polarization
\cite[Lemma~19.2]{MilneAV}. Put $s=\pi+V$. The polynomial
identity for $\pi$ gives
\[
                              s^2=5.
 \tag{RME.2}
\]
It remains to justify integrality of $(s-1)/2$, including the
connected part of the two-torsion group scheme.

The three Weierstrass points over $0,1,\infty$ are rational.
Their three distinct pairwise differences give all nonzero
geometric two-torsion points, since the curve is ordinary of
genus two. Thus $A[2]_{\mathrm{red}}$ is the constant group
scheme $(\mathbf Z/2)^2$. Over the perfect field $\mathbf F_2$,
the reduced subscheme of a finite group scheme is a subgroup;
its map to the etale quotient is an isomorphism. In the
commutative case this splits the group scheme into its connected
and reduced factors. Cartier self-duality of $A[2]$, supplied
by the principal polarization, then identifies its connected
factor with the dual of $(\mathbf Z/2)^2$. Equivalently, the
dual split decomposition contains the connected subgroup
$\mu_2^2$ of rank four, which is already the full connected
rank. Consequently, over the stated field,
\[
                       A[2]\simeq(\mathbf Z/2)^2\times\mu_2^2.
 \tag{RME.3}
\]
Here Cartier duality is at the group-scheme level, not just on
geometric points; see \cite[Theorem~11.1 and the opening of
Section~16, pp.~18 and~29]{MilneAV}.

On the constant etale factor, $\pi$ is the identity, while on
$\mu_2^2$ it is zero, since Frobenius squares its coordinates.
Under Cartier duality and the principal polarization, $V$ has
the opposite two restrictions: zero on the etale factor and
identity on the multiplicative factor. Thus $s$ acts identically
on the entire group scheme $A[2]$.

The endomorphism $s-1$ therefore kills the kernel of $[2]$ as
a finite group scheme. By the quotient property of the isogeny
$[2]:A\to A$, it factors uniquely as $[2]T$ for an integral
endomorphism $T$ defined over $\mathbf F_2$. Substituting
$s=2T+1$ in (RME.2) and using torsion-freeness of the
endomorphism group gives $T^2+T-1=0$. Also $s^\dagger=s$, so
$T^\dagger=T$. As $X^2+X-1$ is irreducible over $\mathbf Q$,
the resulting unital action of its quadratic integer ring is
injective. This is the required integral $\mathcal O_5$ action.
\end{proof}

\paragraph{Its purely two-primary theta candidates.}
For this same curve the Raynaud theta divisor has exactly three
geometric two-primary points, namely the classes
$\mathcal O(W_i)\otimes B^{-1}$ of its three Weierstrass points.
Indeed A24 proves the exact relation
$\Phi^2-\Phi+4=0$ for $\Phi=\pi^2$, the $4$-power Frobenius.
On the geometric etale Tate module, $\Phi$ is invertible, so
this relation gives $\Phi=1$ modulo four. Applying A35 over
$\mathbf F_4$ places every two-primary theta point in
$A(\mathbf F_4)$. The verified count
$\#C(\mathbf F_4)=3$ leaves only those three Weierstrass
points in the Abel image. Their theta classes are distinct
two-torsion points because $2W_i\sim2B\sim K_C$; they are
nonzero since $B$ is non-effective on an ordinary curve.
In particular a non-Weierstrass point of this curve has a
nonzero odd-primary component in its canonical class. This
does not prescribe the prime divisors of that component or
prove that any one of them has even order of two.
